% ============================================================
%  capitulo6_conclusiones.tex
%  Capítulo 6: Conclusiones y Recomendaciones
% ============================================================

\chapter{Conclusiones y Recomendaciones}
\label{chap:conclusiones}

Este capítulo cierra la monografía con las conclusiones principales
organizadas en torno a los cuatro objetivos específicos y el objetivo
general, las recomendaciones dirigidas a los actores clave y las líneas
de investigación futura que este trabajo deja abiertas.


% ─────────────────────────────────────────────────────────────
\section{Conclusiones}
\label{sec:conclusiones}
% ─────────────────────────────────────────────────────────────

\subsection{Sobre el diagnóstico de los barrios subnormales}
\label{subsec:conc_diagnostico}

La magnitud del problema de la subnormalidad eléctrica en Colombia es
mayor de lo que indican los indicadores de cobertura convencionales. Con
más de 314.000 familias en esa condición, el 90\% concentradas en la
región Caribe, y un crecimiento del 26\% en las viviendas sin título
entre 2020 y 2024, la brecha entre la demanda de normalización y los
recursos disponibles del PRONE se amplía año tras año
\parencite{SSPD2020, DANE2024a}.

Un hallazgo que tiene consecuencias técnicas directas: los usuarios de
BS en pisos térmicos cálidos consumen entre el 61\% y el 98\% más de lo
que proyectan los estándares del PIEC \parencite{SSPD2025, UPME2025}.
Diseñar un sistema AGPE con los valores del PIEC en lugar de los datos
reales de la plataforma O3 significa subdimensionarlo, y un sistema
subdimensionado no genera el ahorro en factura que se le promete al
usuario.

La consulta a 20 expertos fue igualmente contundente: el principal
obstáculo para la normalización no es técnico. Es el miedo al incremento
de la factura y la desconfianza acumulada hacia el OR. Esos dos factores
solos explican la mayoría de los proyectos cancelados, bloqueados o
revertidos en el país, y ninguno de los dos se resuelve desde la
ingeniería. Los 20 expertos calificaron el acompañamiento social con
5 sobre 5, sin excepción.

\subsection{Sobre el marco normativo y los casos de estudio}
\label{subsec:conc_normativo}

Entre 2014 y 2025, la regulación colombiana pasó de prohibir
prácticamente la autogeneración a habilitarla dentro de los proyectos
de normalización financiados con recursos públicos del PRONE, contemplar
la cesión de activos a Comunidades Energéticas y permitir la entrega
de excedentes al SIN bajo la CREG 174 de 2021
\parencite{Ley1715_2014, Decreto1023_2024, ResolucionCREG174_2021}.
Es un cambio normativo estructural.

El problema es que ese cambio llegó sin los instrumentos metodológicos
que los OR necesitan para aprovecharlo. La vigencia 2024 lo dejó claro
de la manera más directa posible: ningún proyecto elegible. La norma
existe; la capacidad de implementación no ha madurado al mismo ritmo.

Los tres casos de estudio analizados llevan a la misma conclusión,
independientemente del OR, el periodo o la región: la calidad técnica
de la infraestructura es necesaria, pero no es suficiente. La gestión
social previa, la articulación interinstitucional y la condonación de
las deudas históricas pesan igual o más sobre los resultados a largo
plazo \parencite{CastilloQuintana1991, MendezFuentes2013,
SerranoRueda2019}.

\subsection{Sobre el análisis técnico y financiero}
\label{subsec:conc_tecnico}

Garantizar el consumo básico de subsistencia se ve limitado por la
potencia máxima que puede declarar un sistema AGPE según la normatividad
vigente, especialmente cuando su ubicación se requiere en el sitio de
consumo. Esto se debe a dos factores: primero, si el sistema es conectado
en los transformadores de BT dimensionados en la normalización, no se
puede garantizar la generación para todos los usuarios atendidos en ese
punto; y segundo, si se optara por la tipología tipo granja, la
disponibilidad del terreno se convierte en un cuello de botella porque
los usuarios del BS difícilmente cuentan con terrenos aptos para ello.

Dadas las limitaciones de potencia disponible en los puntos de conexión
y espacio en los sitios de consumo, se propone la autogeneración remota
como la alternativa técnica más viable. No obstante, esto requiere
articulación institucional con el sector público o privado, en el
entendido de que los terrenos tendrían que ser aportados al Estado por la
persona natural o jurídica que pretenda darlos en uso y goce del
proyecto, mínimo hasta la finalización de su vida útil. En las
observaciones realizadas por la ciudadanía al proyecto de resolución de
la convocatoria 001, se evidencia que otra de las alternativas viables
para garantizar este consumo básico es a través de sistemas
individuales; sin embargo, hasta la fecha en las convocatorias se ha
limitado esta tipología.

Para que el esquema de sostenibilidad sea eficiente cuando el AOM lo
pretenda asumir una Comunidad Energética, el análisis de alternativas
debe tener en cuenta que la tecnología seleccionada sea lo menos
compleja posible para que la población pueda participar en su
mantenimiento preventivo, sin que este sea el factor preponderante
frente al potencial energético, el LCOE o la afectación ambiental.

Se observó que para garantizar la sostenibilidad del sistema de AGPE se
debe lograr mínimo una generación de 153~kWh/mes por usuario, ya que
como se evidenció, si este valor es inferior el usuario seguirá viendo
una reducción en la tarifa, pero asociada al esquema de subsidios
creados en la Ley 142 \parencite{Ley142_1994}.

La viabilidad de los proyectos de normalización con AGPE en la región
Caribe no es uniforme: depende de la combinación específica de
condiciones de cada barrio, incluyendo el recurso solar disponible,
la capacidad del transformador más cercano, el número de usuarios y su
consumo real. Estos resultados confirman que la integración de AGPE en
proyectos de normalización requiere un análisis caso a caso, lo que
fundamenta los criterios técnicos y financieros incorporados en las
Fases 2 y 3 de la metodología propuesta.

\subsection{Sobre la metodología propuesta}
\label{subsec:conc_metodologia}

La metodología en cuatro fases secuenciales y una transversal propuesta
en el Capítulo~\ref{chap:metodologia} es la primera herramienta de este
tipo desarrollada en Colombia que articula de manera explícita los
componentes técnicos, normativos, financieros, sociales e
institucionales de los proyectos de normalización con AGPE dentro del
marco del PRONE. Su estructura responde directamente a la brecha que
explica por qué ningún proyecto fue elegible en 2024.

Su mayor aportación no está en el diseño eléctrico — eso está bien
cubierto por el RETIE y los estándares de los OR — sino en la
articulación de los componentes que históricamente han determinado la
diferencia entre el éxito y el fracaso: la gestión social, la
certificación de los BS, la coordinación interinstitucional y la
sostenibilidad de los activos en el largo plazo.

La fase transversal de gestión social y comunitaria es la más crítica.
Los 20 expertos, la literatura y los tres casos de estudio coinciden
en eso. Un proceso de gestión social serio no cabe en una jornada de
información de una hora antes de arrancar las obras. Requiere comenzar
con el diagnóstico del territorio, intensificarse durante las obras y
continuar activo durante al menos el primer año post-normalización. El
presupuesto del proyecto tiene que reflejarlo.

Los proyectos con incidencia directa sobre los usuarios residenciales
de estratos bajos, que están relacionados con la prestación de servicios
públicos domiciliarios y que podrían generar una modificación en la
facturación manejada históricamente, requieren un acompañamiento social
fuerte que garantice el éxito y el funcionamiento a lo largo de la vida
útil de la infraestructura. Así mismo, requieren la articulación con
entidades externas a los prestadores de servicio y del Estado.

Los sistemas de medición individuales y las funcionalidades de la
medición avanzada, entre ellas la prepago, no solo permiten mejorar la
gestión del consumo de energía, sino que tienen una connotación social
ya que inciden en la percepción del usuario sobre los hábitos de consumo
y el uso eficiente y gestión de la energía.

La región Caribe es crítica en los procesos de normalización,
independientemente de si es el Estado o el OR quien aborda la
problemática. Primero, porque es el área donde más se presenta este
fenómeno, y segundo porque el consumo energético es más elevado, sumado
a las condiciones de pobreza y bajo nivel educativo.


% ─────────────────────────────────────────────────────────────
\section{Recomendaciones}
\label{sec:recomendaciones}
% ─────────────────────────────────────────────────────────────

\subsection{Para el Ministerio de Minas y Energía}
\label{subsec:rec_mme}

\begin{enumerate}[label=\textbf{\arabic*.}, leftmargin=1.8cm, itemsep=8pt]

    \item \textbf{Publicar una guía metodológica oficial para la
    formulación de proyectos PRONE con AGPE.} El cero proyectos elegibles
    de 2024 confirma que la necesidad es real. El MME podría adoptar
    o adaptar la metodología propuesta en esta investigación como
    herramienta de soporte para los OR.

    \item \textbf{Incorporar criterios de priorización específicos para
    proyectos con AGPE.} Las convocatorias actuales aplican criterios
    genéricos que no reconocen la mayor complejidad y el mayor costo
    unitario de estos proyectos. Los criterios diferenciados deberían
    valorar el recurso solar disponible, el impacto esperado en la
    factura del usuario, la solidez del plan de AOM y la constitución
    de la Comunidad Energética.

    \item \textbf{Establecer seguimiento plurianual obligatorio.} Exigir
    informes de impacto a los 24 y 48 meses de la energización, con
    indicadores de recaudo, calidad del servicio, estado del AGPE y
    satisfacción comunitaria. Vincular el desempeño en proyectos
    anteriores a la elegibilidad para convocatorias futuras.

    \item \textbf{Definir un esquema de remuneración para las Comunidades
    Energéticas.} En coordinación con la CREG, crear un mecanismo que
    retribuya a la comunidad por la gestión de los activos AGPE cedidos
    con recursos del PRONE, haciendo esa responsabilidad económicamente
    sostenible.

    \item \textbf{Incluir como variable de priorización la socialización
    previa con la comunidad.} Se recomienda que una de las variables de
    priorización sea un documento remitido por el OR donde se evidencie
    que las comunidades beneficiarias conocen de qué trata el proyecto y
    el impacto tarifario que esto conlleva.

    \item \textbf{No exigir la inclusión obligatoria de AGPE en pisos
    térmicos templado y frío.} Se recomienda que en los pisos térmicos
    templado y frío no sea de obligatoria inclusión los sistemas de AGPE
    en los proyectos que presenten los OR en esas áreas de influencia,
    dada la menor disponibilidad del recurso solar.
\end{enumerate}

\subsection{Para los Operadores de Red}
\label{subsec:rec_or}

\begin{enumerate}[label=\textbf{\arabic*.}, leftmargin=1.8cm, itemsep=8pt]

    \item \textbf{Iniciar la gestión de certificaciones con seis meses
    de anticipación.} Sin la certificación del alcalde, no hay proyecto.
    Ese trámite puede tardar meses y depende de voluntad institucional
    que el OR no controla. Empezar tarde es el error más común y el
    más evitable.

    \item \textbf{Construir capacidad interna de formulación.} La brecha
    de 2024 no se resuelve solo contratando consultores. Los equipos del
    OR necesitan entender los requisitos del programa y poder estructurar
    proyectos elegibles con sus propios recursos. Eso requiere equipos
    multidisciplinarios que combinen ingeniería eléctrica, trabajo social
    y gestión financiera.

    \item \textbf{Presupuestar la gestión social en serio.} La experiencia
    de Popayán muestra que ese componente representa una fracción
    significativa del costo total. Tratarlo como gasto accesorio a
    minimizar es una de las razones por las que los proyectos bien
    diseñados técnicamente terminan fracasando en campo.

    \item \textbf{Mostrar la factura post-normalización desde la primera
    jornada de socialización.} Con números reales, con y sin AGPE,
    explicando el subsidio FOES. La incertidumbre sobre cuánto va a
    costar el servicio es el principal generador de resistencia
    comunitaria, y la mejor forma de reducirla es con información
    transparente desde el inicio.

    \item \textbf{Garantizar la concordancia de los documentos
    presentados en las convocatorias.} Se evidenció en las evaluaciones
    publicadas por el MME que se presentaban muchos errores de
    digitación, información incompleta, modificación de los formatos
    exigidos por la entidad, doble contabilidad en los presupuestos,
    inconsistencias en la cantidad de usuarios respecto a la
    certificación de la alcaldía y esquemas de sostenibilidad mal
    planteados.

    \item \textbf{Realizar socializaciones y capacitaciones con la
    comunidad en el primer año posterior a las obras.} Los OR, el MME
    y los entes territoriales deben reforzar la orientación en temas
    como medidas de eficiencia pasivas y activas, uso y gestión eficiente
    de la energía, cadena de valor, costo de la energía según la hora de
    consumo, interpretación de los datos del medidor, esquema de
    subsidios, entre otros.
\end{enumerate}

\subsection{Para las Alcaldías Municipales}
\label{subsec:rec_alcaldias}

\begin{enumerate}[label=\textbf{\arabic*.}, leftmargin=1.8cm, itemsep=8pt]

    \item \textbf{Mantener actualizadas las certificaciones de BS.} Es
    una obligación legal, pero también es la condición que permite a las
    familias acceder al subsidio FOES al que tienen derecho. No
    actualizarlas tiene un costo real sobre las comunidades más
    vulnerables del municipio.

    \item \textbf{Participar activamente en la gestión social de los
    proyectos.} La alcaldía no es un actor que solo firma papeles. Su
    presencia en las jornadas de socialización, su capacidad de gestión
    de conflictos comunitarios y su articulación con otras entidades del
    Estado fueron factores determinantes en Popayán
    \parencite{SerranoRueda2019}. Sin esa participación, los proyectos
    más complejos no se pueden sostener.
\end{enumerate}


% ─────────────────────────────────────────────────────────────
\section{Propuestas de Investigación Futura}
\label{sec:investigacion_futura}
% ─────────────────────────────────────────────────────────────

Este trabajo deja abiertas las siguientes líneas de investigación que
complementarían sus resultados:

\textbf{1. Distribución de pérdidas por nivel de tensión en BS de la
región Caribe.} La regulación vigente no desagrega con precisión las
pérdidas técnicas y no técnicas por nivel de tensión dentro de los BS.
Un estudio que desarrolle esa desagregación permitiría intervenciones
más focalizadas y mejores criterios de priorización del PRONE.

\textbf{2. Impacto de las Comunidades Energéticas en la sostenibilidad
del AGPE.} El Decreto 2236 de 2023 es reciente y no hay experiencias
documentadas de su aplicación en BS colombianos. Un seguimiento de tres
años a la constitución, operación y mantenimiento de una Comunidad
Energética en un BS normalizado generaría información de alto valor para
perfeccionar la regulación y los mecanismos de soporte.

\textbf{3. Interacción entre subsidios energéticos y AGPE en comunidades
de bajos ingresos.} No existe hoy ninguna disposición que articule el
FOES con la generación solar de los usuarios de BS. Un modelo que
evalúe el efecto de distintos esquemas de subsidio sobre la factura de
usuarios normalizados con AGPE podría fundamentar recomendaciones
concretas de política para el MME y la CREG.

\textbf{4. Modelo de optimización para priorizar BS a normalizar con
AGPE.} Los recursos del PRONE son limitados y la demanda supera la
capacidad de intervención. Un modelo que combine variables técnicas
(recurso solar, estado de la red, número de usuarios), sociales
(pobreza, capacidad de pago, presencia étnica) y financieras (costo
por usuario, recaudo proyectado) para priorizar los BS con mayor
potencial de impacto y sostenibilidad sería una herramienta directamente
útil para el CAPRONE.

\textbf{5. Modificación de los esquemas de subsidios existentes.} Se
propone analizar la posibilidad de que los esquemas de subsidios
tengan en cuenta la inversión realizada por el Estado a través de
proyectos financiados por fondos y programas, incluido el PRONE.

\textbf{6. Remuneración vía tarifa para el AOM de sistemas AGPE.}
Se requieren propuestas en torno a la remuneración vía tarifa para los
prestadores de servicio público que realicen la actividad de AOM de los
sistemas de AGPE financiados con recursos públicos y privados, ya que
actualmente no se encuentra regulado por la CREG.

\textbf{7. Regulación de la autogeneración remota.} Análisis y posible
modificación sobre el Proyecto de Resolución 701\_92 de 2025 expedido
por la CREG, relacionada con la regulación de los sistemas de
autogeneración remota.

\textbf{8. Plan de Normalización de Redes Eléctricas a nivel nacional.}
Si bien actualmente el MME cuenta con el Plan Nacional de
Electrificación Rural que asocia acciones sobre el PRONE, y la UPME
realiza Planes Indicativos de Expansión de la Cobertura, no está claro
cómo se aborda integralmente la subnormalidad. Dado que este es un
fenómeno que, por las condiciones socioeconómicas del país, el
crecimiento poblacional y los conflictos armados, no se va a reducir si
no se realiza una intervención planificada y focalizada.